\section{Excursus on Method}

This is a post we had no intention to write and it was not on our schedule. Nevertheless, it seems necessary at this time to summarize several months of efforts. First of all, one must resist the urge to debate, since our goal is gnosis, wisdom, or knowledge. Hurling platitudes across a barrier will never lead to knowledge; that should be verifiable in your own experience. This is not to deny the value of informed discussions that ideally lead to clarification, precision and depth. The difference is that a discussion requires shared presuppositions and a shared goal. I've been told that in Internet lingo there is a creature called a ``troll'', whose only purpose is to disrupt such discussions. However, in traditional terms, a guest who arrives uninvited and then feels it his right to set the tone and agenda is merely rude and vulgar.

Any discussion about whether some spiritual movement is traditional or not must be based on specific criteria, some of which will be listed below. Please try to apply these criteria to all such discussions.

It is important not to be fooled by words without a specific reference to a concrete time and place. For example, as a word ``Christianity'' means very little; on the other hand, it is fruitful to discuss Western Europe from the time of Charlemagne until the Reformation. Note, too, that Guenon went to Egypt not to Albania; there is a reason for that.

\paragraph{Hierarchy}


\begin{itemize}


\item Is the society hierarchical?

\item Is it caste-based (spiritual, warrior, skilled and servile workers)

\item What is the position on spiritual authority?

\item How is temporal power exercised?
\end{itemize}


\paragraph{The Elite and Initiation}


\begin{itemize}


\item Is there an elite who knows and understands metaphysical principles?

\item Is there an initiatic path to lead suitable candidates to such knowledge?

\item Is the foundation of thought based on transcendence or simply material concerns?

\item Are there dogmas, symbols and rites that the lower castes can accept that are consistent with the knowledge of the elite?
\end{itemize}


\paragraph{Cycles}


\begin{itemize}


\item Is there an awareness of understanding of cosmic cycles?

\item Is there knowledge of a Primordial state and a fall from it?

\item How is the decline and ultimate repeat of the cycle understood?
\end{itemize}


\paragraph{Ownness}


\begin{itemize}


\item Does the spiritual system define the society?

\item Does it recognize boundaries (e..g, saved/unsaved, clean/unclean, etc.)?

\item Are such boundaries preserved and protected?
\end{itemize}


To give an example, it is quite clear that in Medieval Europe, the spiritual system was felt and experienced as ``one's own''. Those who opposed it, whether internal or external, were treated as enemies of the state. Nowadays, no one fears excommunication from the church. However, at that time, that it was much more serious since it also implied excommunication from society and the consequent loss of all civil rights.

\paragraph{Decisionism}


In a hierarchical society, each layer pledges loyalty to those above him, and then ultimately to the Emperor who embodies both spiritual authority and political authority in the same man. As such, his spiritual authority is infallible and his political authority is absolute. In making law, the Emperor has no theoretical restraints. Evola wrote about this in \emph{Pagan Imperialism} prior to (I believe) Carl Schmitt's writings on decisionism, which amounts to the same thing.

Therefore, when, for example, an Emperor Constantine decides to change the state religion, he is acting within his rights and his loyal subjects will go along. Similarly, over time, the various kings made the same decision. Those who see an injustice in this reveal their latent modernism. Evola did not seem willing to recognize this as a logical consequence of this view.


\flrightit{Posted on 2011-01-20 by Cologero}
